\section{Problem Formulation}
\label{sec:problem}

\subsection{System description}

The simulated system consists of:
\begin{itemize}
  \item One (single-player) or two (two-player) KUKA LBR iiwa~7 robots,
        each a 7-DOF revolute manipulator with joint state
        $\jq\in\R^{7}$, $\jqdot\in\R^{7}$ and joint limits
        $\jq^{\mathrm{lo}}_{i}\le\jq_{i}\le\jq^{\mathrm{hi}}_{i}$.
  \item A rigid flat-disk paddle attached to the end-effector via a fixed
        constraint at offset $\mathbf{p}_{\mathrm{pe}}=[0,0,0.05]^{\!\top}$
        in the EE frame.
  \item A spherical ball with state
        $\bball(t),\,\vball(t)\in\R^{3}$.
        Between contacts it obeys gravity-only projectile motion
        \begin{equation}
          \ddot{\bball}(t)=-g\,\mathbf{e}_{z},\quad g=9.81\,\mathrm{m/s^{2}},
          \label{eq:gravity}
        \end{equation}
        with small linear damping in the simulator.
  \item A regulation-size table whose surface lies at
        $z=z_{T}=0.76\,\mathrm{m}$, with a net of height
        $z_{N}=0.9125\,\mathrm{m}$ at $x=0$.
\end{itemize}

The world frame is right-handed with the table running along the $x$-axis.
Robot~1 (Player~0) is based at $\mathbf{p}_{1}=(-2.0,0,0.73)$ with zero
yaw and defends $x<0$. Robot~2 (Player~1) is based at
$\mathbf{p}_{2}=(+2.0,0,0.73)$ rotated by yaw $\pi$ about $z$ and defends
$x>0$.

\subsection{Observation and action spaces}

At every $1/240$\,s simulator step the agent receives the dictionary
\[
  \mathbf{o}=\bigl(\bball,\vball,\jq,\jqdot,\;\text{table\_info},\;\text{player\_side}\bigr),
\]
where ball state and own joint state are the relevant signals for control.
The agent emits a target configuration $\jq^{\mathrm{cmd}}\in\R^{7}$.
The simulator's joint controller is in \emph{position mode} with a
\textsc{maxVelocity} cap of $12$\,rad/s per joint and a
$250$\,N$\,$m torque cap, so the realised joint trajectory is the
position-controlled response to a sequence of setpoints rather than a
torque or velocity profile.

\subsection{Rules and termination}

A rally proceeds under the rules below, encoded in the provided
\texttt{SimulationManager}:
\begin{itemize}
  \item The receiving side must return the ball over the net such that it
        bounces on the opponent's half exactly once before the opponent
        contacts it.
  \item Faults that end the rally:
        \emph{double bounce} on a player's own side,
        \emph{ball hits the net},
        \emph{ball goes out of bounds},
        \emph{ball drops below} $z=0.5\,\mathrm{m}$,
        or \emph{rally timeout} ($10\,\mathrm{s}$).
  \item Scoring: a game ends at $11$ points with a $2$-point lead.
  \item In single-player mode the ball is launched from
        $\mathbf{p}_{0}^{\mathrm{serve}}=(+0.6,0,1.1)$ with nominal
        velocity $(-2.8,0,+1.0)$ plus Gaussian noise
        $\sigma_{xy}=0.3$, $\sigma_{z}=0.15$. The receiver's only goal
        is the \emph{successful\_return} fault.
\end{itemize}

\subsection{Algorithm-design problem}

Let $\jq^{\mathrm{cmd}}_{t}\in\R^{7}$ be the joint setpoint commanded at
step $t$. The robot's controller produces actual joint trajectories
$\jq(t)$ that drive the paddle pose
$\mathbf{x}_{p}(t)\in SE(3)$ through the forward-kinematics map
$\mathbf{x}_{p}=\Phi(\jq)$. The ball state propagates under~\eqref{eq:gravity}
between contacts and according to a restitution model when it touches the
table or paddle. We seek a control policy
\begin{equation}
  \pi:\mathbf{o}\mapsto\jq^{\mathrm{cmd}},
  \label{eq:policy}
\end{equation}
that maximises the expected number of \emph{successful\_return} events per
rally (or, in two-player mode, the expected number of games won) subject
to joint, workspace and net-clearance constraints. The policy must be
deterministic from $\mathbf{o}$, may not query the opponent's joint state,
and may not modify the simulator.

The optimisation is hard for three reasons. First, the ball's actual
trajectory differs from idealised projectile prediction because the
PyBullet ball has linear damping and the table coefficient of restitution
differs slightly from the nominal value, so the model is biased. Second,
the controller is position-controlled with a velocity cap, so the
realised Cartesian path is \emph{not} a free design variable --- it is the
forward image of a joint-space linear interpolation. Third, the
two-robot setup is geometrically symmetric, but the manipulator's
inverse kinematics has multiple branches; naïvely sending the right-hand
robot the mirror world target produces a different joint solution that
yields a different (and worse) Cartesian swing path.
